\chapter{Concluding remarks}\label{ch:discussion}

% ⚠ RENAMED AND EXTENDED 2026-09-27, after Abram compared the chapter against the SoftLang template
%   (Chapter "Concluding remarks" = Summary · Limitations or threats to validity · Future work).
%   Chapter 7 is now titled "Concluding remarks" and §7.1 is "Summary", so Chapter 7 IS the
%   template's concluding chapter and the three sections match it one to one. This also answers the
%   gap left when Chapter 8's "Summary and outlook" was dropped the same day -- the thesis had no
%   summary anywhere for about two hours.
%   Labels UNCHANGED on purpose: ch:discussion is referenced by §6.4's closing sentence and
%   sec:disc-findings by nothing outside this file; a key is not prose (CLAUDE.md §6).
% ⚠ MANIFEST §10's "Do not restate the verdicts here" IS SUPERSEDED, not broken. That instruction was
%   written for a Discussion section sitting between the RQ chapters and the limitations. The
%   template's Summary brief is "summarize this work in an insightful manner, assuming that the
%   reader has seen the rest", and a summary that withholds the findings is not one. §10 rewritten
%   with this write.
% THREE SUMMARY PARAGRAPHS ADDED in front of the three lead-ins, 302 words: what was built and asked ·
%   what retrieval returned · what generation returned. No hypothesis ID and no RQ number anywhere --
%   the summary assumes the reader has seen §5.6 and §6.4, so IDs would make it a second pass over
%   them. Every figure is one those two sections already print.
%   Abram's two edits on the draft, taken with the write: the generator size the generation numbers
%   come from is stated, together with the 3-billion replication (§6.4's own nuance, that what fails
%   is narrower than the verdict sounds); and the drafted closing pair -- "None of this survived a
%   larger generator unchanged. Every comparison kept its direction at three billion parameters and
%   two of the three reversed at seven billion." -- was cut as off-register for a summary and
%   replaced by the mechanism, that a larger generator holds more of the answer itself.
%   The old opening's first and third sentences went with the change: the summary now does the work
%   they did, and only the transition sentence before the lead-ins survives.
% Chapter 7 section 7.1: the discussion.  Written 2026-09-27, drafted in the chat as prose over
%   three passes and approved before this write (CLAUDE.md §2 steps 5-6).
% Shape mirrors §5.6 and §6.4: three bold lead-ins, ~450 words, no float.  It carries ONLY what
%   needs both RQs at once (manifest §10) -- no verdict is restated, no hypothesis is named, and
%   claim C-4 stays in §5.3.
% Title is "Discussion", Abram's choice 2026-09-27, replacing the drafted "What the two stages say
%   together" and the stub's "Findings".  Label sec:disc-findings kept -- a key is not prose.
% Every figure is a registry number already printed in Chapter 5 or 6: containment leading on 70
%   then 30 topics (R7); faithfulness 0.189 vs 0.171 and factual correctness 0.206 vs 0.241 (R9);
%   the two concordances 96/105 and 83/105 (R9.B); blocked 0.551 against dense 0.557, 0.006 apart
%   (R1.2, R3.6); SLT-TYPE 0.104 retrieval and 0.049 faithfulness (R3.5, R9).  Nothing new.
% ⚠ C-7 AS WORDED IN THE REGISTRY IS SUPERSEDED BY THE SECOND LEAD-IN.  The registry said a
%   retrieval gain reliable at the retrieval stage does not survive into generation; what the
%   numbers support is the two-direction claim written here -- the ordering over the fifteen reaches
%   attribution tightly and correctness loosely, AND a difference the retrieval measure cannot
%   resolve separates the answers anyway.  Manifest §3 C-7 reworded with this write.
% ⚠ "NOT DRAGGED BELOW" IS DELIBERATE and it is the honest form.  A per-topic check of the 1.5B
%   per-sample score files (2026-09-27, DIAGNOSTIC ONLY -- not in the registry, not reported, not
%   citable) shows SLT-TYPE's 0.196 against the floor's 0.170 is carried by five topics rescued
%   from zero factual correctness, with a median per-topic difference of 0.000.  So "above the
%   floor" overstates it, and "because the generator has effectively fallen back on its own
%   knowledge" explains a five-topic tail as though it were a level.  §6.3.2 and §6.4 carry the
%   matching one-clause repair.
% Cut on Abram's instruction with the write: the two closing sentences of the third lead-in, which
%   named the relevance criteria of §3.1 and prescribed that retrieval for community mathematical
%   question answering deliver formula and prose together.  The third lead-in was retitled from
%   "What a context for mathematical question answering has to carry" to match what remains -- say
%   the word and it goes back.
% The limitations section is named IN WORDS, not by \ref: sec:disc-limitations is still commented
%   out below.  Convert to \ref when §7.2 lands (§6.3 did the same while §6.4 had no label).
\section{Summary}\label{sec:disc-findings}

This thesis built a retrieval-augmented system for mathematical question answering and asked what
structure-aware formula retrieval contributes to it. Three indexes over formula trees were built
beside a sparse and a dense text index, and the rankings they produce were merged in ten
arrangements. Those arrangements were scored on answer retrieval over the 78 topics of the ARQMath-3
collection, under the labels of a relevance judge validated before it was relied upon. The same
fifteen configurations then supplied contexts to a generator, and its answers were scored for how
much of each came from its context and how much of each was correct.

Formula evidence did not improve answer retrieval. Each formula index added to dense retrieval left
the system behind dense retrieval alone, and all three added together left it further behind still.
How the rankings were merged did decide the outcome. Collapsing the three formula rankings into one
vote before fusing beat letting each of them vote separately, 0.551 against 0.374, and that
arrangement is the best multi-index retriever this thesis produced. It does not beat text-only
retrieval. At 0.551 against dense retrieval's 0.557 the two are indistinguishable on both reported
measures.

A retrieved context did reach the answer. The multi-index context made answers more factually
correct than generation without retrieval, 0.206 against 0.170, and both faithfulness and factual
correctness rose with retrieval quality across the fifteen configurations. Against the strongest
text-only retrieval the same context was drawn on more heavily and produced less correct answers:
0.189 faithfulness against 0.171, and 0.206 factual correctness against 0.241. Both readings are
taken at 1.5 billion parameters, and both hold at 3 billion. The gains thin as the generator grows:
at 7 billion parameters a generator holds more of the answer itself, so a context has more to clear
before it helps, and some of these comparisons run the other way.

Three things follow from the two chapters read together, and none of them can be stated inside
either one.

\textbf{Formula overlap misleads a ranker and a generator alike.}
Section~\ref{sec:res-transfer} located why formula evidence does not help answer retrieval. A
topic's formulas are contained in the answers the assessors judged relevant more than in those they
judged not relevant, leading on 70 of the 78 topics; among the best-matching answers, where a
ranking is decided, that lead falls to 30. Formula matching finds the right neighbourhood and
cannot order inside it. Chapter~\ref{ch:generation} shows the generator doing the same thing to the
same posts. Under the multi-index retriever's context more of the answer is traceable to that
context than under the strongest text-only retrieval's, 0.189 against 0.171, and less of what the
answer says is correct, 0.206 against 0.241. The ranking treats formula overlap as pertinence and
places those posts in front of the generator; the generator treats them as pertinent and draws on
them. One signal misleads two stages in the same direction, and it reaches the second stage through
the first.

\textbf{What a retrieval score predicts about an answer.} Across the fifteen configurations of
Section~\ref{sec:met-retrieval-configs} the ordering by the primary measure and the ordering by
faithfulness are concordant on 96 of the 105 pairs, and the ordering by factual correctness on 83.
A retrieval score predicts how much of an answer will come from its context more closely than it
predicts whether that answer will be right. Where the measure resolves nothing, the answers still
differ: the blocked system and dense retrieval are 0.006 apart on the primary measure, inside the
0.02 of Section~\ref{sec:met-judge-validation}, and their answers separate on both generation
measures, in opposite directions. The relation also does not hold evenly at the weak end. SLT-TYPE
retrieves worst of the fifteen, at 0.104, and its answers are not dragged below generation without
a retrieved context, while five configurations that retrieve better than it are. Almost nothing of
its answers comes from its context, at 0.049 faithfulness, and a context that is not drawn on
cannot carry an answer wrong. What decides the cost of a weak context is how heavily it is used,
not how weak it is.

\textbf{What formula evidence alone amounts to.} Section~\ref{sec:ret-verdicts} found that blocking
earns its place by keeping formula evidence from doing harm rather than by supplying what the text
rankings lack. Chapter~\ref{ch:generation} prices the evidence that still reaches the generator. The
three configurations built from formula indexes alone are among the five whose answers fall below
generation without a retrieved context. A well-matched formula, delivered without the prose that
resolves the question, is not a usable context at either stage: the ranking cannot order inside the
neighbourhood it finds, and the generator writes a worse answer from it than from nothing at all.

Each of these readings rests on conditions measured once, on one collection, and scored by
instruments this thesis validated rather than assumed. Section~\ref{sec:disc-limitations} states what that
leaves unsettled.

% Chapter 7 section 7.2: the limitations. Written 2026-09-27, rewritten the same day on Abram's
%   reading of the first version: it was over-polished -- long sentences, abstract vocabulary,
%   semicolon chains -- where the house voice is short declaratives carrying one fact each, with a
%   longer sentence connecting them (CLAUDE.md §3). Snapshots `chapter_07_discussion.pre72.bak.tex`
%   (the stub) and `chapter_07_discussion.pre72Simple.bak.tex` (the over-polished version).
% FORM, all three on Abram's instruction 2026-09-27:
%   (i) a NUMBERED LIST, not one subsection per threat -- superseding manifest §10's "one subsection
%       per threat, each opening with a bold Threat to validity: <level> (<type>) label". tocdepth is
%       2, so each subsection would spend a table-of-contents line on a three-sentence item, and
%       limitations pass the house itemize test (swap any two, no inference breaks).
%   (ii) NO validity type in brackets after the bold name. All five are construct threats; the label
%       repeated five times carried nothing.
%   (iii) NO opening paragraph -- the section goes straight into the list. What went with it: the
%       sentence drawing the boundary against the guards ("the guards of §4.3, §4.4 and §5.3 defended
%       the assumptions that evidence could defend; this section states the threats that remain"),
%       which was the one place the thesis stated that boundary in prose, and the sentence saying all
%       five are construct threats. Manifest §10's guard/limitation boundary is now kept only by
%       where things sit, not by any sentence.
% FIVE ITEMS, Abram's selection: the relevance judge · the generation reference · no human reading ·
%   factual correctness and the length of an answer · answering efficiency. The last four are his own
%   notes, polished.
% ⚠⚠ THE POST-HOC DISCLOSURE IS NOT HERE, AND THIS WAS ITS ONLY SITE. Manifest §10's post-hoc row
%   requires §7.2 to disclose that H1.3 was restated on 2026-09-20 with R3 in view, that H3.1's idea
%   came from the retired ladder (and that the observation does not reproduce), and that H3.1's
%   clause (a) was restated after its data in the direction that made it easier to pass. §5.6 carries
%   no disclosure (Abram, 2026-09-21) and §6.4's was moved here (Abram, 2026-09-26). With §7.2
%   written without it, THE THESIS DISCLOSES NONE OF THE THREE. Raised at the write, left out on his
%   instruction, manifest §10 reopened. Add before submission.
% ⚠ ALSO DROPPED on his instruction: what the retrieval measure scores against what Chapter 6 asks of
%   it (which would have discharged the reservation flagged in `_sec64_rq2_verdicts.tex` -- still
%   undischarged); one collection and one test year; the stack as built rather than optimised;
%   choices stated as method and never varied; one run per condition with no statistical test. So
%   §7.2 concedes no EXTERNAL and no INTERNAL threat at all.
% The three scope decisions (no statistical test, no comparison with published answer-retrieval
%   systems, retrieval reported under the judge's labels alone) are stated NOWHERE in Chapter 1 --
%   it has no scope section. Manifest §10 expects them in §1.6. Nothing here cross-references them.
% Numbers, all read from the live sources: Cohen's kappa 0.174 and the 38 official participant runs
%   (§4.4); the 0.02 fuzziness value (§4.4, §5.1); up to five judged answers per topic (§6.2).
% κ is referenced to §4.4, where its value for this judge is reported, and not also to §3.6, which
%   defines it -- one pointer, per CLAUDE.md's localised-term rule.
% Trimmed 2026-09-27 on Abram's reading (snapshot `chapter_07_discussion.pre72Trim.bak.tex`):
%   (a) "the measure can therefore mark a correct answer down, and it cannot mark a wrong one up" ->
%       "a missing route can only cost a correct answer; it cannot help a wrong one". Same one-way
%       bias, said plainly; the old wording read as two operations on a score rather than one gap.
%       Note the claim is deliberately about the MISSING ROUTE, not about the measure as a whole --
%       a wrong answer with correct intermediate claims can still score well, so "it cannot make a
%       wrong answer look right" would have been too strong.
%   (b) "and the proposal this work began from asked about it" CUT -- it pointed the reader at a
%       document they have not seen. ⚠ The efficiency item therefore no longer says WHY the omission
%       is owed, so manifest §10's "concede it rather than dropping the word silently" is now met
%       only by the item existing. Its last sentence was rebuilt to stand alone: it says nothing here
%       shows how the system behaves under a query load, rather than leaving "that question" hanging.
%   (c) the closing sentence handing over to §7.3 CUT. §7.2 now ends on the list, so the chapter has
%       no sentence carrying the reader from the limitations into the future work -- against
%       CLAUDE.md §3's "end a section by naming the problem the next one inherits". His call.
% Never here: a guard (those are §4.3, §4.4 and §5.3), a definition, or a result.
\section{Limitations}\label{sec:disc-limitations}

\begin{enumerate}[leftmargin=1.7em,itemsep=0.7em,topsep=0.5em]

\item \textbf{The relevance judge.} Every retrieval score in Chapter~\ref{ch:retrieval} is computed
from labels a language model assigned. Section~\ref{sec:met-judge-validation} shows that the judge
orders systems as the assessors do, which is what this thesis reads it for, and that it is lenient
about individual labels, agreeing with them at a Cohen's $\kappa$ of 0.174. The scores therefore say
which system is better, not how good any of them is. Two scores differing by 0.02 or less say nothing
at all, and that limit was measured on the 38 official participant runs of one year. Whether it holds
on another collection is untested.

\item \textbf{The generation reference.} Factual correctness is measured against the union of up to
five judged answers for each topic. A mathematics question usually has more than one route to the
result. Those five answers may take several of those routes or may all take the same one, and either
way they do not cover every route. An answer that is right by a route none of them takes counts as
unsupported.

\item \textbf{No human reading of the generated answers.} Both generation measures are computed by
the evaluation framework's own language model. Section~\ref{sec:met-judge-validation} validates the
relevance judge, and Section~\ref{sec:gen-method} says that this is a different instrument doing a
different task. Nothing here shows that these factual-correctness scores agree with a mathematician
reading the same answers. Nothing here judges an answer as a whole either. Both measures work over
the separate claims an answer is split into, so no result in Chapter~\ref{ch:generation} says that an
answer is right or that it is wrong.

\item \textbf{Factual correctness and the length of an answer.} Factual-correctness precision is the
share of an answer's claims that the reference supports. A longer answer makes more claims, and each
further claim is one the reference may not cover. A low score can therefore come from length rather
than from error. Nothing holds the answers to a common length, and the length they run to differs from
one condition to the next. Some of what separates two conditions may be how much each answer said.

\item \textbf{Answering efficiency was never measured.} How fast a system answers is part of whether
it can be used. Every run reported here was made offline, in batch. No latency was timed and no query
throughput was counted, so nothing here says how the system would behave under a query load.

\end{enumerate}

% Chapter 7 section 7.3: future work. Written 2026-09-27, drafted in the chat and approved with
%   edits before this write (CLAUDE.md §2 steps 5-6). Snapshot `chapter_07_discussion.pre73.bak.tex`.
% THREE ITEMS, Abram's own (he said "two of them" and listed three; all three were drafted and he
%   kept all three): prompt tuning · which formula to query with · one representation for a formula
%   and its prose. Numbered list, no opening paragraph -- the same form as §7.2, so the chapter reads
%   as one piece.
% Each item names the finding it follows from, so none of it is a wish list: item 1 from the single
%   prompt of §6.2 and the 0.225 faithfulness ceiling (§6.4, registry R9.A); item 2 from fanout (§4.3,
%   §5.2.2) and §5.5's erosion at the narrow end; item 3 from §7.1's third reading.
% "ATTENTION" IS USED BARE, on Abram's instruction 2026-09-27 -- "use attention word only, it is easy
%   and direct". ⚠ Chapter 3 does not define it and there is no definition site anywhere, so this is
%   a term used before it is defined, which is the defect the first-draft review named. Defensible
%   here because §7.3 measures nothing and names a direction rather than a method; flagged in case an
%   examiner stops on it. The drafted paraphrase that avoided the word was "a model that weighed each
%   formula in a question against the words around it".
% CUT on his instruction with the write, both closing sentences:
%   (i) item 2's "Section 5.5's diagnostic is the measure to test it against: the separation between
%       relevant and non-relevant answers holds over a topic's whole judged set and erodes among the
%       best-matching answers, which is where a ranking is decided." So item 2 proposes the direction
%       without naming how it would be evaluated.
%   (ii) item 3's "The learned systems of Section 2.2 already take that shape and lead at formula
%       search, and nothing here tests whether they carry to answer retrieval." So §7.3 cites no
%       related work and SSEmb is not named here -- the honest concession against it now rests on
%       §2.2 alone, since the §7.2 item that conceded no retriever training was also dropped.
% The chapter now ends on "would not need the merge at all", which is what Chapter 8 picks up.
\section{Future work}\label{sec:disc-future-work}

\begin{enumerate}[leftmargin=1.7em,itemsep=0.7em,topsep=0.5em]

\item \textbf{Tuning the answer-generation prompt.} This thesis presents one prompt and holds it
fixed, so that what separates its conditions is the context and nothing else. That leaves the prompt
itself unexamined. The highest faithfulness anywhere in Chapter~\ref{ch:generation} is 0.225, so most
of every answer comes from the generator rather than from what was retrieved. A prompt that directs
the generator more firmly at its context would raise that share, and whether it also raises factual
correctness is the question to put to it.

\item \textbf{Choosing which formula to query with.} The formula indexes here are queried with every
non-trivial formula a topic contains. Section~\ref{sec:res-transfer} shows what that costs. With so
many query formulas, almost every plausible answer matches one of them, and the ranking has little
left to order them by. An attention mechanism over the question, trained to score each formula by how
much of the question it carries, would pick the one or two that matter and query with those.

\item \textbf{One representation for a formula and its prose.} Every configuration in this thesis
keeps the two apart. Text rankings and formula rankings are produced separately and then merged, and
Section~\ref{sec:disc-findings} gives the reason that arrangement can only go so far: a formula match
locates the right neighbourhood, and the words around the formula decide which answer in it resolves
the question. A retriever trained to embed a post as one object, with its formulas and its prose in a
single representation, would not need the merge at all.

\end{enumerate}
